\section{Introduction}

Deep reinforcement learning (RL) has matured into a workbench for studying
machine-learning methodology that does not appear in supervised pipelines:
non-stationary data distributions, sparse credit assignment over thousands of
time steps, and the tight coupling between exploration policy and the very
data the model trains on. These properties make RL a uniquely useful
pedagogical setting for a graduate deep-learning course --- the same
optimizer, the same backbone, and the same loss family produce radically
different outcomes once a policy is allowed to choose its own training data.

We study these effects on \textbf{Slime Volleyball}~\citep{slimevolleygym}, a
two-player continuous-physics game that compresses the methodology questions
into a tractable form. The state is 12-dimensional (positions and velocities
of agent, ball, and opponent), the action space is \textsc{Discrete}(6), and
each episode lasts up to 3000 frames or until one side loses all five lives.
The environment is small enough that a full ablation matrix fits in a
weekend's GPU budget, and large enough that subtle algorithmic choices
produce statistically meaningful differences. We benchmark every variant
through (a) a deterministic-greedy reward against the scripted baseline and
(b) a 1500-match round-robin Elo ladder against all other variants, and we
deploy each trained policy as an ONNX model in a browser playable so that the
ablation is reproducible end-to-end.

Our contributions are:
\begin{itemize}
\item A clean $2\times2\times2$ ablation matrix over Dueling, Double DQN, and
  $n$-step returns on a shared trunk (Section~\ref{sec:results}) that
  isolates each Rainbow~\citep{hessel2018rainbow} component's
  contribution and exposes a destructive interaction: $n$-step returns
  \emph{without} Double DQN are worse than vanilla DQN.
\item A self-play scaffold that combines a FIFO opponent
  pool~\citep{silver2016alphago,vinyals2019alphastar} with a
  baseline-mixing knob, plus an empirical demonstration that pure-pool
  self-play \emph{underperforms} mixing the scripted baseline as an
  occasional opponent.
\item A PPO~\citep{schulman2017ppo} entropy-collapse case study with a
  four-change rescue intervention (entropy coefficient $10\times$,
  rollout length doubled, linear LR anneal, target-KL early stop) that
  swings final reward by $+3.6$.
\item An ``ecosystem effect'' finding: a long-standing $\times 10$ scale
  bug in the in-pool baseline produced a lineage of trained agents that
  dominate the public ladder (Elo $\sim$1800), and a correctly-trained
  agent loses to them at $\sim$8:1 odds in head-to-head play.
\item Browser-deployable ONNX exports of every variant (one shared
  inference contract: 12-dim float input, six logits / Q-values out)
  with $<10^{-5}$ max-abs verification against the PyTorch reference
  forward pass.
\end{itemize}

The remainder of the paper formalizes the MDP (Section~\ref{sec:problem}),
specifies architectures, losses, and self-play training
(Section~\ref{sec:methods}), reports ablation results
(Section~\ref{sec:results}), discusses safety and robustness implications
(Section~\ref{sec:safety}), and concludes (Section~\ref{sec:conclusion}).
